\section{Results and Discussion}
\label{sec:results}
\subsection{Operating reference for the flexibility assessment}
The selected day's coordinated trajectory establishes what the facility would be doing before receiving an additional grid request. Figure~\ref{fig:day} shows the corresponding grid import and workload execution. Computation moves between eligible intervals, with lower import around some price peaks and higher import when work is completed or storage is charged. The arrival profile remains unchanged: scheduling alters execution, not the computational demand presented to the facility.

\begin{figure*}[t]
\centering
\subfloat[Grid import and electricity reference price.]{\includegraphics[width=0.49\textwidth]{Figure_4a.png}}\hfill
\subfloat[Workload arrival and execution.]{\includegraphics[width=0.49\textwidth]{Figure_4b.png}}
\caption{Reference and coordinated operation on 22 September 2025. The figures' ``baseline'' denotes reference operation; the coordinated trajectory is the undisturbed reference for the subsequent event assessment. Opening states are inherited from the preceding day.}
\label{fig:day}
\end{figure*}

The scheduling pattern reflects more than the instantaneous price. Work already deferred has a shrinking execution window; charging is limited by power ratings and available capacity; and changes in IT execution alter cooling requirements. Thus neither installed capacity nor the local price alone determines the response still available at a given time. The heatmap and component results below assess that remaining capability from the actual coordinated trajectory.

\subsection{Magnitude, duration and direction of response}
Figure~\ref{fig:envelope} reports certified feasible duration for each requested deviation and start time. Reduced-import requests of 100 and 150 kW are sustained for six hours from noon. This corresponds to reductions of 10\% and 15\% of rated IT power, while the actual target is defined at the whole-facility meter. A 25 kW increased-import request reaches the twelve-hour assessment cap at 03:00 and at 10:00, 11:00 and 12:00. These cap-reaching results establish twelve-hour feasibility under the recovery conditions; they do not establish an intrinsic upper limit of twelve hours.

\begin{figure*}[t]
\centering\includegraphics[width=\textwidth]{Figure_6.png}
\caption{Certified feasible duration on 22 September 2025. Negative deviations reduce import; positive deviations increase it. Blank cells mean that no positive duration was certified by the search. A plus sign marks an unresolved next 15-minute trial. Durations reaching twelve hours or the remaining local-day length are assessment-capped. Other unmarked positive cells have an adjacent tested boundary; the search does not establish global maximality for every cell.}
\label{fig:envelope}
\end{figure*}

The dominant feature is the dependence on operating time. A 100 kW reduction lasts 1.25 h from 06:00 and 6 h from noon, despite unchanged equipment ratings. The storage state, unfinished workload and thermal trajectory differ between these starts. Those conditions determine both what can be curtailed during the event and whether its consequences can be recovered afterwards. The absence of a certified response at another time should therefore not be interpreted as the absence of installed flexible equipment.

Within the tested grid, larger requested magnitudes generally have shorter certified durations, with plateaus where different magnitudes reach the same duration. The two directions also differ. Small import increases can be assembled over long periods by executing eligible work earlier than in the coordinated schedule and by changing storage and cooling operation. Import reduction depends on being able to defer execution, reduce chiller demand or supply IT from the UPS. Since the tested increased-import requests are substantially smaller than the reduced-import requests, the comparison does not isolate direction from magnitude. It demonstrates the capabilities evaluated here, rather than a complete symmetric power range.

Two assessment boundaries are important when reading the figure. First, an event cannot extend beyond the selected local day, although recovery may continue for three hours afterwards. The declining late-day durations consequently reflect this assessment limit as well as physical capability. Second, thirteen cells retain a plus sign because the next candidate duration remained unresolved. For example, the 25 kW increase at 06:00 is feasible for at least 11.25 h, and the 50 kW increase for at least 8.75 h. A solver timeout cannot establish that a longer response is impossible.

The reported durations include the requirement to recover relative to undisturbed operation. This prevents a response from appearing feasible merely because it leaves storage depleted or additional work outstanding beyond the comparison window. Nevertheless, the heatmap remains a state-conditioned illustration from one day. It does not indicate how frequently these durations would be available during the year, or the dependable capacity that could be offered into a particular market.

\subsection{How internal resources deliver the grid response}
The component results explain why a single facility-level deviation can conceal very different operating actions. Figures~\ref{fig:reduction} and \ref{fig:increase} show changes from the coordinated trajectory, with positive and negative component changes summed according to Equation~\eqref{eq:decomposition}. Their interpretation is electrical: TES is represented through changes in chiller electricity, and UPS through net charging minus discharging. The allocation is one accepted cost-minimising dispatch under the request constraints, rather than a unique division of the response among assets.

\begin{figure*}[t]
\centering\includegraphics[width=\textwidth]{Figure_7.png}
\caption{Component changes for selected reduced-import events. The dashed line identifies the requested deviation during the event; the target is not imposed in the grey post-event region. Recovery is verified by terminal constraints, rather than fully displayed in these panels.}
\label{fig:reduction}
\end{figure*}
\begin{figure*}[t]
\centering\includegraphics[width=\textwidth]{Figure_8.png}
\caption{Component changes for selected increased-import events. Opposing component actions can be larger than the net requested deviation. Grey regions, where present, follow the certified event; the target is released during recovery.}
\label{fig:increase}
\end{figure*}

The reduced-import cases provide a particularly clear comparison. At 06:00, the selected 100 and 200 kW reductions are delivered entirely through additional UPS discharge for 1.25 h; IT and chiller demand remain at the coordinated values. At 15:00, both requested reductions are sustained for 3 h, but the response is supplied mainly by deferred IT execution and reduced direct cooling. TES is full and idle in the undisturbed trajectory at that start, so there is no TES charging to curtail. The 200 kW case is delivered by IT and cooling alone; the 100 kW case includes only a single 37 kW UPS-discharge interval. Installed energy is therefore not a reliable guide to which resource will participate in an event.

The IT--cooling interaction is visible in the afternoon cases. Moving computation changes both its direct electrical demand and the timing of heat generation. Cooling consequently contributes to the metered reduction alongside IT. A workload-only power model would omit that thermal consequence, while a cooling model driven by an unchanged heat profile would miss the new operating opportunity. The present results demonstrate these coupled mechanisms. They do not quantify a cost or flexibility premium over separately optimised subsystems, which would require an explicit comparison with such schedules.

Increased-import events show another form of coordination. At midnight, the selected 25 and 75 kW responses contain substantial internal changes that partly cancel at the meter. At 06:00, work already eligible for execution can be processed earlier than in the undisturbed schedule, with associated changes in cooling and storage. This permits the 25 kW increase to continue for at least 11.25 h and the 75 kW increase for 7.75 h. Advancing execution here means using an earlier eligible interval than the coordinated schedule selected; it never means processing a job before it arrives.

Long-duration provision therefore need not correspond to one device sustaining a constant change throughout the event. The facility can transfer the burden between resources as workload obligations and physical states evolve. One component may even move against the requested direction while another compensates. That behaviour is essential to interpreting the large opposing bars in Figure~\ref{fig:increase}: only their sum is seen at the grid connection. The stored results reconcile component sums and target tracking within the specified numerical tolerances (\supp{7}).

Together, these examples show what integration contributes to flexibility characterisation. It identifies which combinations are feasible from a given state, explains why the participating resources change with time, and includes the subsequent recovery obligation. A single quoted flexibility capacity would conceal all three features.

\subsection{Annual electricity-cost value}
Over the complete 2025 trajectory, reference operation costs GBP~550,503.59 and coordinated operation GBP~522,702.95. The resulting saving is GBP~27,800.64, or 5.050\%. These totals compare the same arriving workload and thermal operating limits; reference cooling is already cost optimised. The result therefore measures the additional electricity-cost benefit of scheduling eligible work and coordinating storage under the observed reference prices.

Annual grid electricity increases from 6.718 to 6.794 GWh, approximately 1.13\%, while peak import rises from 974.9 to 1,681.9 kW. Lower expenditure is consequently not evidence of lower energy consumption or a smoother load profile. The objective rewards the timing of electricity use. Conversion losses and changes in utilisation can increase energy use while the shift towards lower-priced intervals reduces cost. Peak charging and workload execution are also admissible when their combined import remains below the modelled structural limit. A real connection agreement or capacity charge could change the attractiveness of that schedule.

\begin{table*}[htbp]
\centering\small
\caption{Annual signed electricity-cost composition (GBP). Changes are coordinated minus reference; displayed values are rounded.}\label{tab:cost-components}
\begin{tabular}{@{}lrrr@{}}\toprule
Component & Reference & Coordinated & Change \\\midrule
IT electrical demand & 407,653 & 392,409 & -15,244 \\
CRAC chiller & 105,337 & 91,252 & -14,084 \\
TES charging chiller & 0 & 6,903 & +6,903 \\
UPS net grid effect & 0 & -5,376 & -5,376 \\
Auxiliary overhead & 37,514 & 37,514 & +0 \\
Total settlement cost & 550,504 & 522,703 & -27,801 \\
\bottomrule\end{tabular}
\end{table*}


Table~\ref{tab:cost-components} separates the signed electricity-cost terms. IT demand and direct CRAC cooling contribute reductions of approximately GBP~15,244 and GBP~14,084, respectively. TES charging costs GBP~6,903, and net UPS exchange reduces cost by GBP~5,376. Auxiliary expenditure is unchanged. These terms reconcile to the whole-facility saving but are not standalone asset profits. The positive TES charging term records the electricity used to produce stored cooling; its benefits appear partly in reduced direct chiller demand and in the opportunities it creates for the rest of the schedule. Similarly, the change in CRAC cost reflects the timing of IT heat as well as cooling decisions.

Using the full year includes both favourable and less favourable price periods, with no independent daily reset of storage or work. This strengthens the temporal basis of the economic result. It does not establish that 2025 represents every price regime, nor that the resulting percentage applies to a complete site electricity bill. On the selected event day prices are non-negative, whereas 672 of the year's 35,040 quarter-hour intervals have negative prices. The annual valuation therefore encompasses operating conditions that are absent from that single physical illustration.

\subsection{Resource sensitivity and the value of additional UPS capacity}
The annual sensitivities show increasing savings as flexible workload share, UPS energy capacity or TES energy capacity is increased over the tested ranges (Figure~\ref{fig:sensitivity}; full values in \supp{6}). Workload share has the largest response: halving it reduces savings to 3.522\%, while increasing it by half raises them to 6.257\%. UPS capacity spans 4.565--5.501\%, and TES capacity 4.800--5.167\%. This ordering applies to the selected parameter changes. It is not a comparison at equal investment cost, and it does not establish how simultaneous changes would interact.
\begin{figure*}[t]
\centering\includegraphics[width=\textwidth]{Figure_9.png}
\caption{Full-year resource sensitivities relative to the central saving, in percentage points. Panel (a) uses the realised flexible share of the daily input CPU-hours after clipping. Note the different vertical scales. Storage power limits remain fixed as energy capacity changes.}
\label{fig:sensitivity}
\end{figure*}

The workload response reflects both the size of IT demand and its link to heat production. Changing eligible workload creates scheduling opportunities throughout the year, whereas extra stored energy remains constrained by charging power, conversion efficiency and opportunities to use it. The smaller TES response is consistent with retaining fixed thermal charge/discharge ratings and a shared chiller limit. Workload-share sensitivity is therefore useful for interpreting a context-dependent assumption; TES sensitivity shows that a larger energy buffer alone does not remove the other constraints on cooling operation.

UPS enlargement poses a more direct design question: how much additional expenditure could its additional operational saving support? Increasing capacity from 600 to 750 kWh reduces annual coordinated cost by GBP~1,292.28. Increasing it to 900 kWh reduces cost by GBP~2,481.41. These increments correspond to GBP~8.62 and GBP~8.27 per added nameplate kWh per year. The second 150 kWh increment is worth less than the first, although the taper is modest. Because the reserve fraction remains 50\%, these additions increase the operational energy window by only 75 and 150 kWh, respectively.

\begin{table}[htbp]
\centering\small
\caption{UPS enlargement screen: constant annual benefit, ten years, 8\% real discount rate and no additional non-electricity costs.}\label{tab:ups-screen}
\begin{tabular}{@{}rrrr@{}}\toprule
Capacity & Added energy & Added saving & Break-even\\(kWh) & (kWh) & (GBP/year) & investment (GBP) \\\midrule
750 & 150 & 1,292.28 & 8,671 \\
900 & 300 & 2,481.41 & 16,650 \\
\bottomrule\end{tabular}
\end{table}


Under the illustrative ten-year, 8\% real-discount assumptions, Table~\ref{tab:ups-screen} gives supported incremental investments of approximately GBP~8,671 and GBP~16,650, equivalent to GBP~57.81 and GBP~55.50 per added nameplate kWh. These are break-even expenditure thresholds, not estimates of installed battery prices. They identify the amount available to cover enlargement if the observed annual savings persisted. Additional installation and operating costs must fit within that amount for the screen to be favourable.

This result gives a bounded interpretation of flexibility-oriented UPS design. The existing energy-arbitrage benefit supplies a contribution towards larger capacity, but the analysis does not establish that enlargement pays for itself. A deployment decision would need compatible installed costs, degradation, replacement and reserve requirements. It would also need to distinguish a larger UPS with the modelled proportional reserve from an expansion that keeps absolute backup energy fixed. Initial and terminal stored energy differ between capacities, so the single-year saving is used as a conditional annuity rather than a demonstrated lifetime cash flow. Alternative discount rates and lives are reported in \supp{8}. No hypothetical payment for the physical events is added to the annual saving.

\subsection{Practical interpretation and limitations}
The flexibility results indicate how coordinated assets can alter demand while preserving modelled computational obligations and recovering relative to undisturbed operation. They are useful inputs to a wider grid-service assessment, but they do not establish network benefit on their own. The annual cost-minimising schedule has a higher import peak, and the model contains neither network power flows nor the behaviour of other loads. A service procurer would also need to specify response speed, event frequency, delivery tolerances and availability. Quarter-hour trajectories cannot demonstrate subsecond response performance.

The facility and workloads are deliberately stylised. The repeated daily workload, fixed ambient temperature and fixed COP exclude seasonal and weekly changes in computing demand and cooling conditions. Aggregate CPU-hours do not represent server placement, task dependencies, memory constraints or every SLA. The conservative cooling rule further defines the available thermal operating envelope. These assumptions allow the interaction of the represented resources to be examined, but the resulting durations should not be transferred directly to a different facility or scaled solely by its IT rating. GPU power management, application-specific performance trade-offs and other forms of flexibility are outside the present workload-shifting model.

Economic and physical conclusions also have different scopes. Annual results quantify signed reference-price electricity costs, with conversion losses but without battery degradation, network tariffs or service revenue. They do not measure the total financial value of flexibility. The event study establishes feasible responses on one selected day, not an annual availability distribution. It includes no preconditioning and requires recovery within a fixed allowance, so its results are conditional on those service rules. Finally, numerical feasibility checks do not constitute experimental validation of the aggregate facility. Solver bounds, duration-search qualifications and existing terminal-treatment checks are reported in the supplement so that these numerical limits remain distinct from the model's physical assumptions.
